\chapter{Durchführung} 
\label{ch:Durchfuehrung}

\section{Messverfahren}
\label{sec:Messverfahren}

\subsection{Messung von Silber}

Für die Messung der Halbwertszeit von aktiviertem Silber wird eine Geiger-Müller-Zählrohrspannung zwischen $500~\mathrm{V}$ und $550~\mathrm{V}$ am Betriebsgerät eingestellt. Zunächst wird die Untergrundzählrate bestimmt, um diese später von den Messwerten abziehen zu können. Dazu wird die Torzeit auf $10~\mathrm{s}$ eingestellt und $50$ Messungen durchgeführt, sodass die Untergrundzählrate über einen Zeitraum von etwa acht Minuten erfasst wird.

Anschließend wird die eigentliche Silbermessung durchgeführt. Die Träger mit den Silberblechen werden zur Aktivierung für sieben Minuten in der Neutronenquelle platziert. Nach Beendigung der Aktivierung wird das Präparat schnellstmöglich an das Zählrohr gebracht und dort mit einem Aluminiumblech fixiert. Die Messung startet mit einer Torzeit von $10~\mathrm{s}$ über eine Gesamtzeit von $400~\mathrm{s}$. Dieser Zyklus wird insgesamt viermal wiederholt, um die statistische Unsicherheit zu reduzieren. Die vier Messreihen werden anschließend addiert und gemeinsam ausgewertet.

\subsection{Messung von Indium}

Für die Messung der Halbwertszeit von aktiviertem Indium wird analog vorgegangen. Das Messintervall wird dabei auf $120~\mathrm{s}$ eingestellt. Das aktivierte Indium-Präparat wird in der Halterung fixiert und über einen Zeitraum von etwa $50$ Minuten gemessen. Auch hier wird zuvor die Untergrundzählrate bestimmt, wobei die Torzeit entsprechend angepasst wird. Da bei der Aktivierung von Indium neben dem Grundzustand auch ein metastabiler Zustand mit sehr kurzer Halbwertszeit gebildet wird, wird der erste Messpunkt für den Fit vernachlässigt, da dieser vom eigentlichen Zerfall der betrachteten Isotope abweicht.

\section{Versuchsaufbau}
\label{sec:Versucshaufbau}

Der Versuchsaufbau besteht aus mehreren Komponenten. Eine Neutronenquelle, bestehend aus einem Gemisch von Berylliumspänen und dem $\alpha$-Strahler Americium-241, wird von einem Paraffinblock umgeben, der die schnellen Neutronen auf thermische Energien abbremst. Die zu aktivierenden Proben (Silber- und Indiumbleche) werden in einer Probenhalterung positioniert und während der Aktivierungsphase in die Nähe der Neutronenquelle gebracht.

Nach der Aktivierung wird das Präparat in die Position eines Geiger-Müller-Zählrohrs überführt, das mit einem Betriebsgerät verbunden ist. Der Impulszähler zählt die registrierten Zerfälle über jeweils festgelegte Zeitintervalle (Torzeiten). Die Daten werden auf einem PC gesammelt und anschließend mit geeigneter Software ausgewertet.

% \subsubsection*{Skizze des Versuchsaufbaus}
% \label{sec:Skizze}